% Fuente: Auxiliar 5, P1.
\textbf{(Teorema de Carathéodory)} Sea $\mathbf{A}_1,\dots,\mathbf{A}_n$ colección de vectores en $\mathbb{R}^m$.

\begin{enumerate}
\item[(a)] Sea
\begin{equation}
C = \left\{ \sum_{i=1}^{n} \lambda_i \mathbf{A}_i \;\middle|\; \lambda_1,\dots,\lambda_n \ge 0 \right\}.
\end{equation}
Demuestre que todo elemento de $C$ puede expresarse en la forma $\sum_{i=1}^{n} \lambda_i \mathbf{A}_i$, con $\lambda_i \ge 0$, y con a lo más $m$ de los coeficientes $\lambda_i$ no nulos. \textit{Hint:} Considere el poliedro
\begin{equation}
\Lambda = \left\{ (\lambda_1,\dots,\lambda_n) \in \mathbb{R}^n \;\middle|\; \sum_{i=1}^{n} \lambda_i \mathbf{A}_i = \mathbf{y},\; \lambda_1,\dots,\lambda_n \ge 0 \right\}.
\end{equation}

\item[(b)] Sea $P$ la envolvente convexa de los vectores $\mathbf{A}_i$:
\begin{equation}
P = \left\{ \sum_{i=1}^{n} \lambda_i \mathbf{A}_i \;\middle|\; \sum_{i=1}^{n} \lambda_i = 1,\; \lambda_1,\dots,\lambda_n \ge 0 \right\}.
\end{equation}
Demuestre que todo elemento de $P$ puede expresarse en la forma $\sum_{i=1}^{n} \lambda_i \mathbf{A}_i$, donde $\sum_{i=1}^{n} \lambda_i = 1$ y $\lambda_i \ge 0$ para todo $i$, con a lo más $m+1$ de los coeficientes $\lambda_i$ no nulos.
\end{enumerate}

